%%% File:      b09_References.tex
%%% Author:    Joaquín Ataz-López et al
%%% Begun:     July 2020
%%% Concluded: August 2020
%%% Contents:  Just like the previous chapter, initially, this
%            chapter was considered to be a section of Chapter 12.  But
%            when I began writing it I saw that it affected the document
%            as a whole, so it changed place.  In this case the
%            information comes fundamentally from the wiki.
%
%%% Edited with: Emacs + AUCTeX - And at times vim + context-plugin
%%%

% TODO: JAL uses the word "label" for what ConTeXt calls a "reference",
% because he thought it was clearer.  But 2/3 the way through here he notes
% that decision betrayed him, because ConTeXt also has a "label".
% Q: should the terminology be changed throughout this chapter to match
% ConTeXt's?

\environment introCTX_env.tex

\startcomponent b09_References.tex

\startchapter
    [title=References and hyperlinks]

\TocChap

\startsection
    [title=Reference types]

Scientific and technical documents abound in references:

\startitemize

  \item Sometimes these references refer to other documents that are the
    basis for what is being said, or that contradict what is being
    explained, or that develop or further nuance the idea being dealt with,
    etc.  In these cases the reference is said to be {\em external} and, if
    the document is to be academically rigorous, the reference takes the
    form of {\em citations} from the literature.

  \item But it is also common for a document, in one of its sections, to
    refer to another of its sections, in which case the reference is said
    to be {\em internal}.  There is also an internal reference when a point
    in the document comments on some aspect of a particular image, table,
    note, or element of a similar nature, referring to it by its number or
    by the page on which it is found.

    For the purposes of precision, internal references need to be aimed at
    an exact and easily identifiable place in the document.  Hence these
    kinds of references are always a reference to either numbered elements
    (as, for example, when we say \quotation{see table 3.2} or
    \quotation{Chapter 7}), or page numbers.  Vague references of the
    \quotation{as we have already said} or \quotation{as we will see
    further on} kind are not true references, and there is no special
    requirement for typesetting them, nor is there any special tool for
    doing so.  Also, I personally dissuade my PhD or MA students from any
    habitual use of this practice.

\startSmallPrint

Internal references are also commonly called \quotation{cross references}
though in this document I will simply use the term
\quotation{references} in general, and \quotation{internal references}
when I wish to be specific.

\stopSmallPrint

\stopitemize

In order to clarify the terminology I am using for references, I will call
the point in the document where a reference is introduced the {\em origin},
and the location to which it points, the {\em target}.  Seen this way, we
would say that a reference is an internal one when the origin and target
are in the same document, and an external one when origin and target are in
different documents.

From the point of view of typesetting the document:

\startitemize

  \item External references pose no special problem and therefore, in
    principle, do not require any tool to introduce them: all the data I
    need from the target document are available to me and I can use them in
    the reference.  However, if the document of origin is an electronic
    document and the target document is also available on the Web, then it
    is possible to include a hyperlink in the reference that allows one to
    jump directly to the target.  In these cases the document of origin can
    be said to be {\em interactive}.

  \item By contrast, internal references do pose a challenge for
    typesetting the document, since anyone who has experience in the
    preparation of moderately long scientific and technical documents knows
    that it is almost inevitable that numbering of pages, sections, images,
    tables, theorems or similar, to which the cross-reference points, will
    change during the document's preparation, which makes it very difficult
    to keep it up to date (if done by hand).

    \startSmallPrint

        In pre-computer times, authors avoided internal references; and
        those that were inevitable, such as the table of contents (which,
        if accompanied by the page number of each section, is an example of
        an internal reference), were written at the end.

    \stopSmallPrint

\stopitemize

Even the most limited typesetting systems, such as word processors,
allow for the inclusion of some kind of internal cross-references such
as tables of contents.  But that is nothing compared to the comprehensive
reference management mechanism included in \ConTeXt, which can also
combine the internal reference management mechanism aimed at keeping
references up to date, with the use of hyperlinks which is obviously not
exclusive to external references.

\stopsection


\startsection
  [
    reference=sec:references,
    title=Internal references,
  ]

Two things are needed to establish an internal reference:

\startitemize[n]

  \item A label or identifier at the target point.  While compiling,
    \ConTeXt\ will associate particular data with this label.  What data
    will be associated depends on the kind of label it is; it can be the
    section number, the note number, the image number, the number
    associated with a particular item in a numbered list, the section
    title, etc.

  \item A command at the point of origin. This accesses the data associated
    with the label linked to the target point, and inserts it at the point
    of origin.  The command varies depending on which data from the label we
    want to insert at the point of origin.

\stopitemize

When we think about a reference, we do so in terms of
\quotation{origin~\longrightarrow~target}, so it might seem that matters
relating to the origin should be explained first, and then those
relating to target. However, I believe that it is easier to understand
the logic of references if the explanation is reversed.

\startsubsection
  [
    reference=sec:target labels,
    title=The label in the reference target,
  ]

In this chapter, by {\em label} I mean a text string that will be
associated with the target point of a reference and used internally to
retrieve certain information regarding the target point of a reference such
as, for example, page number, section number etc.  In fact, the information
associated with each label depends on the procedure for creating
it.  \ConTeXt\ calls these labels {\em references}, but I think that this
latter term, as it has a much broader meaning, is less clear.

The label associated with the target reference:

\startitemize

  \item Needs each potential target in the document to be unique, so that
    it can be identified without ambiguity.  If we use the same label for
    different targets, \ConTeXt\ will not throw a compiling error but it
    will cause all references to point to the first label it finds (in the
    source file) and this will have the side effect that some of our
    references may be wrong, and, worse still, that we do not notice them.
    Therefore, it is important to make sure, when creating a label, that
    the new label we are assigning has not already been assigned before.

% Careful reformatting this next paragraph, we want extra spaces in the
% second quotation.

  \item Can contain letters, digits, punctuation marks, blank spaces,
    etc.  Where there happen to be blank spaces, \ConTeXt's general rules
    regarding these kinds of characters still apply (see
    \in{Section}[sec:spaces]), so that, for example,
    \quotation{\tt My nice label} and \quotation{\tt My \ \ \
        \ nice \ \ \ \ label} are seen as the same, even
    though they use a different number of blank spaces.

\stopitemize

Since there is no limitation as to which or how many characters can be part
of the label, my advice is to use label names that are clear, and will help
us to understand the source file when, perhaps, we read it long after it
was originally written.  That's why the example I gave before
(\MyKey{My nice label}) is not a good example, as it does not
tell us anything about the target the label is pointing to.  As an example,
for this section the label \MyKey{sec:Target labels} would be
much better.

To associate a particular target with a label there are basically two
procedures:

\startitemize[n]

  \item By means of an argument or command option used to create the
    element to which the label will point.  From this point of view, all the
    commands that create some kind of structure or text element open to
    being a reference target include an option called \MyKey{reference}
    that is used to include the label.  Occasionally, in place of the {\em
    option} the label is the content of the whole argument.

    We find a good example of what I am trying to say in the section
    commands that, as we know from (\in{Section}[sec:sectionsyntax]), allow
    for several kinds of syntax.  In the classic syntax the command is
    written as:

    \type{\section[Label]{Title}}

    and in the syntax specific to \ConTeXt\ the command is written as

    \starttyping
\startsection
  [title=Title, reference=Label, ... ]
    \stoptyping

    In both cases the command provides for the introduction of a label that
    will be associated with the section (or chapter, subsection, etc.\null)
    in question.

    % TODO: The Spanish just has "lists", not "numbered lists".  Should the
    % following paragraph be modified?

    I've said that this possibility is found in {\em all commands} that
    allow us to create a text element that can be the target of a
    reference.  These are all text elements that can be numbered, including
    among others, sections, floating objects of all kinds (tables, images
    and similar), footnotes or end notes, quotations, numbered lists,
    descriptions, definitions, etc.

    \startSmallPrint

        When the label is entered directly with an argument, and not as an
        option to which a value is assigned, it is possible with \ConTeXt\
        to associate several labels with a single target. For example:

        \type{\chapter[label1, label2, label3] {My chapter}}

        I am not clear what the advantage could be to have a number of
        different labels for the one target and suspect \Conjecture that it
        can be done not because it offers advantages but due to some {\em
        internal} requirement of \ConTeXt\ applicable to certain kinds
        of arguments.

    \stopSmallPrint

  \item By means of the \PlaceMacro{pagereference}\tex{pagereference},
    \PlaceMacro{reference}\tex{reference}, or
    \PlaceMacro{textreference}\tex{textreference} commands, whose syntax is:

    \starttyping
\pagereference[Label]
\reference[Label]{Text}
\textreference[Label]{Text}
    \stoptyping

    \startitemize

      \item The label created with \tex{pagereference} allows us to
        retrieve the page number.

      \item Labels created with \tex{reference} and \tex{textreference}
        allow us to retrieve the page number as well as any associated text
        that is included as an argument.

        In both \tex{reference} and \tex{textreference}, the text that is
        linked to the label does not appear in the final document at the
        point where the command is located (the reference target), but can
        be retrieved and printed at the point of origin of the reference.

    \stopitemize

\stopitemize

I said earlier that each label is associated with certain information
regarding the target point.  What that information is depends on the type
of label it is:

% TODO: should remember be in quotes, rather than italicized?

\startitemize

  \item All labels {\em remember} (in the sense that they make it possible
    to retrieve) the page number of the command that created them.  For
    labels attached to sections that may have several pages, that number
    will be the page number where the section in question begins.

  \item Labels inserted with the command that creates a numbered text
    element (section, note, table, image, etc.\null) {\em remember} the number
    associated with that element (section number, note number, etc.)

  \item If this element has a {\em title}, as happens for sections, but
    also tables (and similar) if they have been inserted using the
    \tex{placetable} command, they will {\em remember} this title.

  \item Labels created with \tex{pagereference} only {\em remember} the
    page number.

  \item Those created with \tex{reference} or \tex{textreference} also
    remember the associated text that these commands take as an
    argument.

    \startSmallPrint

        In fact I am not sure of the real difference between
        the \tex{reference} and \tex{textreference} commands.  I think it
        is possible \Conjecture that the design of the three commands that
        allow the creation of labels attempts to run parallel with the
        three commands that allow the retrieval of information from the
        labels (which we will see in a moment); but the truth is that,
        according to my tests, \tex{reference} and \tex{textreference} seem
        to be redundant commands.

    \stopSmallPrint

\stopitemize

\stopsubsection

% TODO: the title of this subsection is crying out for a simplification.

\startsubsection
    [title=Commands at the reference point of origin for retrieving data
        from the target point]

The commands that I will explain next retrieve information from the labels
and, in addition, if our document is interactive, generate a connection to
the reference target.  But the important thing about these commands is the
information that is retrieved from the label.  If we only want to generate
the connection, without retrieving any information from the label, we must
use the \tex{goto} command explained in \in{Section} [sec:createlinks].

\startsubsubsection
    [title=Basic commands for retrieving information from a label]

Bearing in mind that each label associated with a target point can store
different items of information, it is logical that \ConTeXt\ includes three
different commands for retrieving such information; depending on which
information from a reference target point we want to retrieve, we choose one
of these commands:

\startitemize

  \item The \tex{at} command allows us to retrieve the label's page number.

  \item For labels that remember an element number (section number, note
    number, item number, table number, etc.\null) in addition to the page
    number, the \tex{in} command allows us to retrieve this number.

  \item Finally, for labels that remember a text associated with a label (a
    section title, an image title inserted with \tex{placefigure}, etc.\null) the
    \tex{about} command allows us to retrieve this text.

\stopitemize

These three \PlaceMacro{at}\tex{at}, \PlaceMacro{in}\tex{in}, and 
\PlaceMacro{about}\tex{about} commands have the same syntax:

\starttyping
\at{Text}[Label]
\in{Text}[Label]
\about{Text}[Label]
\stoptyping

\startitemize

  \item \MyKey{Label} is the label from which we want to retrieve information.

  \item \MyKey{Text} is the text written just before the information we
    want to retrieve with the command.  Between the text and the data of
    the label that the command retrieves, a non-breakable space will be
    inserted, and if the interactivity function is enabled in such a way
    that the command, besides retrieving the information, generates a link
    that allows us to jump to the target point, the text included as an
    argument will be part of the link (it will be clickable text).

\stopitemize

So, in the following example we see how \tex{in} retrieves the section
number and \tex{at} the page number.

\startDoubleExample
\starttyping
In \in{Section}[sec:target labels], 
that begins on \at{page}
[sec:target labels], the
characteristics of labels used for
internal references are explained.
\stoptyping

In \in{Section}[sec:target labels], that begins on \at{page} [sec:target
labels], the characteristics of labels used for internal references are
explained.

\stopDoubleExample

Note that \ConTeXt\ has automatically created hyperlinks (see
\in{Section}[sec:interactivity]), and that the text taken as an argument
by \tex{in} and \tex{at} is part of the link.  But had we written it
otherwise, the result would be:

\startDoubleExample
\starttyping
In section \in{}[sec:target
labels], that begins on page \at{}
[sec:target labels], the
characteristics of labels used for
internal references are explained.
\stoptyping

% TODO: explain what \about does.


In section \in{}[sec:target labels], that begins on page \at{}
[sec:target labels], the characteristics of labels used for internal
references are explained.

\stopDoubleExample

The text remains the same, but the words {\em section} and {\em page}
that precede the reference are not included in the link as they are no
longer part of the command.

If \ConTeXt\ is unable to find the label that the \tex{at}, \tex{in} or
\tex{about} commands point to, no compiling error will result, but where
the information retrieved by these commands should appear in the final
document we will see \quotation{??} written.

\startSmallPrint

    There are two reasons why \ConTeXt\ cannot find a label:

    \startitemize[n]

      \item We made a mistake when writing the label (either at the origin
        or the target).

      \item We are compiling only a part of the document, and the label
        points to the part not yet compiled (see \in{Sections}[input] and
        \in{}[sec-projects]).

    \stopitemize

    In the first case the error will need to be fixed.  Therefore, it is a
    good idea when we finish compiling the complete document (and the
    second case is no longer an issue), to look for all the appearances of
    \quotation{??} in the PDF to check that there are no {\em broken}
    references in the document.

\stopSmallPrint

\stopsubsubsection


\startsubsubsection
    [title=Retrieving information associated with a label with
        the \tex{ref} command]
\PlaceMacro{ref}

Each of \tex{at}, \tex{in} and \tex{about} retrieve a specific element of a
label.  Another command is available that allows us to retrieve any element
of a given label.  This is the \tex{ref} command, whose syntax is:

\type{\ref[Element to retrieve][Label]}

where the first argument can be:

\startitemize

  \item {\tt text}: returns the text associated with a label.

  \item {\tt title}: returns the title associated with a label.

  \item {\tt number}: returns the number linked to a label.  For example,
    in sections, the section number.

  \item {\tt page}: returns the page number.

  \item {\tt realpage}: returns the actual page number.

  \item {\tt default}: returns what \ConTeXt\ considers to be the {\em
    natural} element of the label.  Generally this coincides with what is
    returned by {\tt number}.

\stopitemize

% TODO: "frontmatter" might be better than "preamble" below.
%        Which of these was used earlier in the book?

In fact, \tex{ref} is much more precise than \tex{at}, \tex{in} or
\tex{about}; for example, it differentiates between the page number and the
actual page number.  The page number may not coincide with the actual page
number if, for example, the page numbering of the document started at 1500
(because this document is the continuation of a previous one), or if the
pages of the preamble were numbered with Roman numerals and the numbering
was restarted after the preamble ended.  Similarly, \tex{ref}
differentiates between the {\em text} and the {\em title} associated with a
reference, something that \tex{about}, for example, does not do.

If \tex{ref} is used to request information from a label which does not
have that particular information (e.g., the title of a label associated
with a footnote), the command will return an empty string.

\stopsubsubsection


\startsubsubsection
    [title=Detecting where the link leads to]

\ConTeXt\ also has two commands that are sensitive to {\em the link
direction}.  By \quotation{link direction}, my intention is to
determine whether the link target in the source file is found before or
after the origin.  For example: we are writing our document and we want to
refer to a section that could still come before or after the one we are
writing in the final table of contents\unknown\ we just haven't decided
yet.  In this situation it would be useful to have a command that writes
one of two texts, depending on whether the target ultimately comes before or
after the origin in the final document.  For needs like this, \ConTeXt\
provides the \PlaceMacro{somewhere}\tex{somewhere} command whose syntax is:

\type{\somewhere{Text if before}{Text if after}[Label]}.

\page[preference]

For example, in the following text, I have used two actual labels in this
chapter in the source file.  One comes before the following text and one
comes after.  (The words \quotation{before} and \quotation{after} are not
red in the following text, not because of the \tex{somewhere} command, but
because they are links, which have been configured to have the specific
colour seen in the example.)

\starttyping
The hyperlink's address can also be detected by the \type{\somewhere} command.
This way we can also find chapters or other text elements
\somewhere {before}{after} [sec:references] and discuss their descriptions
in some other place \somewhere{before}{after} [sec:interactivity].
\stoptyping

\startnarrower\switchtobodyfont[small]

    \color[red]{The hyperlink's address can also be detected by the
     \type{\somewhere} command.  This way we can find chapters or other text
     elements \somewhere {before}{after}[sec:references] and discuss their
     descriptions in some other place \somewhere {before}{after}
     [sec:interactivity].}

\stopnarrower

Another command capable of detecting whether the label it points to
comes before or after is \PlaceMacro{atpage}\tex{atpage}, whose syntax
is:

\type{\atpage[label]}

This command is quite similar to the previous one, but instead of
allowing us to write the text ourselves, depending on whether the label
comes before or after, \tex{atpage} inserts a default text for each of
the two cases and, if the document is interactive, also inserts a
hyperlink.

% TODO: I think the \atpage explanation would really benefit from an example.

The text that \tex{atpage} inserts is the one associated with the
\MyKey{precedingpage} labels in case the {\tt label} it takes as an
argument is {\em before} the command, and \MyKey{hereafter} in the
opposite case.

\startSmallPrint

    When I arrived at this point, I was betrayed by a previous decision: in
    this chapter I decided to call what \ConTeXt\ calls a
    \quotation{reference}, a \quotation{label}.  It seemed clearer to me.
    But certain text fragments generated by \ConTeXt\ commands, such as
    \tex{atpage}, are also called \quotation{labels} (this time in another
    sense).  (See \in{Section}[sec:labels]). I hope the reader does not get
    confused.  I think the context lets us properly distinguish which of the
    different meanings of {\em label} I am referring to in each case.

\stopSmallPrint

Therefore, we can change the text inserted by \tex{atpage} in the same
way that we change the text of any other label:

\starttyping
\setuplabeltext[en][precedingpage=New text ]
\setuplabeltext[en][hereafter=New text ]
\stoptyping

\startSmallPrint

% TODO: This needs to be reviewed for LMTX.

    On this point I believe there is a small error in \suite- (the
    distribution I am using).  Examining the names of the predefined labels
    in \ConTeXt\ that can be changed with \tex{setuplabeltext} there are
    two pairs of labels that are candidates to be used by \tex{atpage}:

    \startitemize[packed]

      \item \MyKey{precedingpage} and \MyKey{followingpage}.
      \item \MyKey{hencefore} and \MyKey{hereafter}.

    \stopitemize

    We could assume that \tex{atpage} would use either the first or the second
    pair.  But in fact, for items coming before, it uses \MyKey{precedingpage}
    and for those following it uses \MyKey{hereafter}, which I think is
    inconsistent.

\stopSmallPrint

\stopsubsubsection

\stopsubsection


% TODO: I can't find what the following refers to.  The URL from ANSS was
% https://wiki.contextgarden.net/References.  Upon Garulfoo's
% reorganization of the wiki, the old URL shows a page which refers to the
% URL below.  But reading that, I don't see anything.
% Can anyone identify what JAL was talking about, and, if it still exists,
% where it is now?

% TODO #2: at the end JAL refers to \showreferences, but when I try it I
% see no output in the PDF nor in the .log or .tuc file.  If this can be
% convinced to do something useful, it should be noted here.
% See https://www.mail-archive.com/ntg-context@ntg.nl/msg78085.html for
% something arguably related, but I can't imagine that is what JAL had in mind.

%% 
%% \startsubsection
%%     [title=Automatic generation of prefixes to avoid duplicate labels]
%% 
%% In a large document it is not always easy to avoid duplication of
%% labels.  It is therefore advisable to put some order into the way we
%% choose which labels to use.  One practice that helps is to use prefixes
%% for the labels that will vary according to the type of label.  For
%% example \MyKey{sec:} for sections, \MyKey{fig:} for figures,
%% \MyKey{tbl:} for tables, etc.
%% 
%% With this in mind, \ConTeXt\ includes a collection of tools that allow:
%% 
%% \startitemize
%% 
%%   \item \ConTeXt\ itself to automatically generate labels for all the
%%     allowable elements.
%% 
%%   \item Every label generated manually to take a prefix, either one we have
%%     predetermined ourselves, or automatically generated by \ConTeXt.
%% 
%% \stopitemize
%% 
%% The detailed explanation of this mechanism is lengthy and, although they
%% are undoubtedly useful tools, I do not think they are essential.
%% Therefore, as they cannot be explained in a few words, I prefer not to
%% explain them and refer to what is said about them in the \ConTeXt\
%% reference manual or in the
%% \goto{wiki}[url(https://wiki.contextgarden.net/References_notes_and_floats/References)] on this
%% matter.
%% 
%% % If we choose to write our own labels, a command that can help us avoid
%% % duplicates is \tex{showreferences}: this command will show a list of
%% % all established labels in the document.
%% 
%% \stopsubsection

\stopsection


\startsection
  [
    reference=sec:interactivity,
    title=Interactive electronic documents,
  ]

Only electronic documents can be interactive; but not all electronic
documents are.  An {\em electronic} document is one that is stored in a
computer file and can be opened and read directly on screen.  On the other
hand, an {\em interactive} electronic document is one that is equipped with
utilities that allow the user {\em to interact} with it; that is: we can do
more than just read it.  There is interactivity, for example, when the
document has buttons that perform some action, or links through which we
can jump to another point in the document (or to an external document), or
when there are areas in the document where the user can write, or there are
videos or audio clips that can be played, etc.

All documents generated by \ConTeXt\ are electronic (since \ConTeXt\
generates a PDF that is by definition an electronic document), but they are
not always interactive.  To provide them with interactivity it is necessary
to expressly indicate this as shown in the next section.

Bear in mind, though, that although \ConTeXt\ can generate interactive PDFs,
in order to appreciate this interactivity we need a PDF reader capable of
it, since not all PDF readers allow us to use hyperlinks, buttons and
similar features typical of interactive documents.

\startsubsection
    [title=Enabling interactivity in documents]
\PlaceMacro{setupinteraction}

\ConTeXt\ does not create a PDF with interactive functions unless expressly
instructed to do so; this is normally done in the preamble of the document.
The command that enables this feature is:

\type{\setupinteraction[state=start]}

Normally this command would be used only once and in the document
preamble when we want to generate an interactive document.  But in fact
we can use it as often as we want by altering the document's
interactivity state. The \MyKey{state=start} command enables
interactivity, while \MyKey{state=stop} disables it, so we can disable
interactivity in some chapters or {\em parts} of our document where we
want to do so.

\startSmallPrint

I can't think of any reason why we would want to have non-interactive parts
in documents that are interactive.  But a key part of \ConTeXt's philosophy
is that if something is technically possible, even if we are unlikely to
use it, a procedure for doing so is provided.  It is this philosophy that
gives \ConTeXt\ so many possibilities, and prevents a simple introduction
like this from being {\em brief}.

\stopSmallPrint

Once interaction is established:

\startitemize

  \item Certain \ConTeXt\ commands include hyperlinks by default.  For example:

    \startitemize

      \item The commands for creating tables of contents will, in
        principle and unless expressly indicated otherwise, be interactive;
        i.e., clicking on an entry in the table of contents will jump to
        the page where the corresponding section begins.

      \item The commands for internal references that we have seen in the
        first part of this chapter provide interactivity: clicking on them
        causes a capable PDF reader to jump to the reference target.

      \item For footnotes and endnotes, a click on the note anchor in the
        main body of the text will take us to the page where the note
        itself is written; further, a click on the note mark in the note
        text will take us to the point in the main text where the reference
        to the footnote or endnote was made.

      \item Etc.

    \stopitemize

  \item Other commands specifically designed for interactive documents,
    such as presentations, are available.  Presentations employ numerous
    tools associated with interactivity such as buttons, menus, image
    overlays, embedded sound or video, and so on.  Explaining all this
    would be too long and besides, presentations are a rather special kind
    of document.  Therefore, in the following lines I will only describe
    one feature associated with interactivity: hyperlinks.

\stopitemize

\stopsubsection


\startsubsection
    [title=Basic configuration for interactivity]

\tex{setupinteraction}, in addition to enabling or disabling interaction,
allows us to configure some matters related to it; mainly, but not only,
the colour and style of the links.  This is done through the following
command options:

\startitemize

  \item {\tt color}: controls the {\em normal} colour of links.

  \item {\tt contrastcolor}: determines the colour of links where the
    target is on the same page as the origin.  I recommend that this option
    always be set to the same content as the previous one.

  \item {\tt style}: controls the link style; for
    example, \quotation{\type{style=\it}} would cause links to be typeset
    in italics.  By default links are set in bold face.

  \item {\tt title, subtitle, author, date, keyword}: The values assigned
    to these options will be converted into metadata of the PDF generated
    by \ConTeXt.

  \item {\tt click}: This option controls whether the link should be
    highlighted when it is clicked.

\stopitemize

\need 7\baselineskip
For example, one day the \tex{setupinteraction} command for this document
was as follows:
\starttyping
\setupinteraction
[
    state=start,
    color=darkblue,
    contrastcolor=darkblue,
    style=,
    title={A Not So Short Introduction to ConTeXt LMTX (Community Edition)},
    author={Joaquín Ataz-López et al},
]
\stoptyping

\stopsubsection

\stopsection


\startsection
    [title=Hyperlinks to external documents]

I will distinguish between commands that do not create the link but help
to typeset the URL of the link, and commands that create the hyperlink.
Let's look at them separately.

\startsubsection
    [title={Commands that help typeset the hyperlinks but do not create them}]

URLs tend to be very long, and include characters of all types, even
characters which are reserved characters in \ConTeXt\, and therefore cannot
be used directly.  Further, while it might be desirable to have a clickable link
comprised of a short \quotation{human-readable} word, in other cases it
might be desired to write out the entire URL\null.  (If the PDF reader
implements facilities for interactive documents, hovering the mouse cursor
over a link should cause the reader to display the link; however, if
someone prints the PDF, the actual target of the link won't be seen.)

In addition, when the URL itself is displayed in the document, it may be
very difficult to typeset it nicely, as the URL may not contain acceptable
places for \ConTeXt\ to insert a line break.  Longer URLs can exceed the
length of a line, but even shorter URLs may by chance come toward the end
of a line in a paragraph, possibly causing line breaking problems.
Moreover, it is not reasonable to hyphenate URLs to insert line breaks, as
the reader could hardly know whether or not the hyphen actually forms part
of the URL\null.

% TODO: JD's tests indicate that \useURL does a pretty good job of line
% breaking with long URLs.  So JAL's text seems a bit outdated.  For
% starters, I modified the next paragraph.

\ConTeXt\ provides two utilities for {\em typesetting} URLs, which will be
discussed in turn.
%
% The first is primarily for URLs that will be used internally, but will not
% actually be displayed in the document.  The second is for URLs that have to
% be written in the text of the document.  Let's look at them separately.
The first does a reasonable job of breaking URLs when necessary.  The
second may be more capable, but the URLs it typesets are not clickable
(i.e., they are not interactive).

\need 5\baselineskip

\startdescription{\tex{useURL}}\PlaceMacro{useURL}

This command allows us to write a URL and associate a name with it, so that
when we want to use the URL in our document we can invoke it with the
associated name.  We can use \tex{useURL} either in the body of the
document at a convenient point (but before we want to use the URL), or in
the preamble of the document.  In the case of URLs that will be used
several times throughout the document, it is especially useful to define
them in the preamble.

\page[preference]

The command implements two usages:

\startitemize[n, packed]

  \item \type{\useURL [Associated name] [URL]}
  \item \type{\useURL [Associated name] [URL] [] [Link text]}

\stopitemize

\startitemize

  \item In the first version, the URL is simply associated with the name by
    which it will be invoked in our document.
    % TODO(?): JAL originally said the following, but I don't understand
    % that, at least vis-a-vis what I see when I run LMTX.
    % But then, to use the URL, we will have to indicate somehow, when
    % invoking it, which clickable text will be shown in the document.
    That is, when the URL is invoked, the full URL is displayed in the
    document, and the entire URL is clickable text.

%  Example of overfull line:
% 
%     \useURL[pppp][http://www.pragma-ade.comgeneralmanualsoioioioioioioiooiowhy-is-there-air.pdf]
% This sis siisisis     \from[pppp]
% 
%     lkslds sdsd dddks d \hyphenatedurl{http://www.pragma-ade.comgeneralmanualsoioioioioioioiooiowhy-is-there-air.pdf}
% 
% \hyphenatedurl{https://amalgamated.consolidated.incorporated.com/topic/subtopic/subsub/and/enough/to/demonstrate/line/breaking/here.html}.  This URL is \quotation{nicely} broken across lines in this paragraph.

  \item In the second version, the last argument provides the clickable
    text to be typeset in the document.  The third argument exists in case
    we want to divide a URL into two parts, so that the first part contains
    the access address and the second part the name of the specific
    document or page that we want to open.  For example, consider the
    address of the document that explains what \ConTeXt\ is:
    \color[blue]{\hyphenatedurl{http://www.pragma-ade.com/general/manuals/what-is-context.pdf}}.
    This address can be written in full in the second argument, leaving the
    third empty:

\starttyping
\useURL [WhatIsCTX]
  [http://www.pragma-ade.com/general/manuals/what-is-context.pdf]
  []
  [What is \ConTeXt?]
\stoptyping

but we can also split it into two arguments:

\starttyping
\useURL [WhatIsCTX]
  [http://www.pragma-ade.com/general/manuals]
  [what-is-context.pdf]
  [What is \ConTeXt?]
\stoptyping

\useURL [WhatIsCTX]
  [http://www.pragma-ade.com/general/manuals]
  [what-is-context.pdf]
  [What is \ConTeXt?]

In both cases we will have associated that address with the word
\MyKey{WhatIsCTX}; to include a link to that address, when we use the link
creation command (as described in \in{Section}[sec:createlinks]), instead
of the URL itself, we simply use \MyKey{WhatIsCTX} as the argument to
that command, giving us the link \quotation{\from[WhatIsCTX]}.  You can see
that this text is in the usual link colour for this document, and if you
are reading this with a PDF reader capable of interaction, when you hover
the mouse over the text, it should show you the target of the link.

If at any point in the text we want to reproduce a URL that we have
associated with a name using \tex{useURL}, we can use the
\tex{url[Associated name]} which inserts the URL associated with that name
into the document.  But this command, although it typesets the URL into the
document, does not create any link, as seen here: \type{\url[WhatIsCTX]}
produces \quotation{\url[WhatIsCTX]}.

\startSmallPrint

    The format in which the URLs written using \tex{url} are displayed is
    not the one established by \tex{setupinteraction}, but rather the one
    specifically established for this command by 
    \PlaceMacro{setupurl}\tex{setupurl}, which allows us to set the style
    (option {\tt style}) and the colour (option {\tt colour}).  In the
    preamble for this document the same colour was specified in both cases,
    which is why the above URL is the same colour as clickable links.

\stopSmallPrint

\stopitemize

\stopdescription

\startdescription{\tex{hyphenatedurl}}\PlaceMacro{hyphenatedurl}

This command is intended for URLs that will be written out in full in the
text of our document.  This command instructs \ConTeXt\ to include line
breaks within the URL, if necessary, to correctly typeset the paragraph.
Its format is:

\type{\hyphenatedurl{URLaddress}}

Despite the name of the \PlaceMacro{hyphenatedurl}\tex{hyphenatedurl}
command, it does not hyphenate the name of the URL\null.  What it does is to
consider that certain characters common in URLs are good points to insert a
line break (either before or after these characters).  \ConTeXt\ allows us
to add the additional characters to the list of characters where a line
break is allowed.  We have three commands for this:

\starttyping
\sethyphenatedurlnormal{Characters}
\sethyphenatedurlbefore{Characters}
\sethyphenatedurlafter{Characters}
\stoptyping\PlaceMacro{sethyphenatedurlnormal}\PlaceMacro{sethyphenatedurlbefore}\PlaceMacro{sethyphenatedurlafter}

These commands add the characters they take as arguments to, respectively,
the list of characters that allow line breaks before and after, the list
of characters that only allow line breaks before, and those that only
allow line breaks after.

% TODO: should we consistently use "an URL" or "a URL" ?

\tex{hyphenatedurl} can be used whenever it is desired that the full text
of a URL should appear in the final document (more or less as is).  It can
even be used as the last argument to \tex{useURL} in the version of that
command where the last argument picks up the clickable text to be displayed
in the final document.  For example:

\need 4\baselineskip
\starttyping
\useURL [WhatIsCTX]
  [http://www.pragma-ade.com/general/manuals/what-is-context.pdf]
  []
  [\hyphenatedurl{http://www.pragma-ade.com/general/manuals/what-is-context.pdf}]
\stoptyping

\need 4\baselineskip

In the \tex{hyphenatedurl} argument, all of \ConTeXt's normally reserved
characters can be used, except three which must be replaced by commands:

\startitemize[packed]

  \item \%{} must be replaced by
    \PlaceMacro{letterpercent}\tex{letterpercent}

  \item \#{} must be replaced by \PlaceMacro{letterhash}\tex{letterhash}

  \item \backslash{} must be replaced by
    \PlaceMacro{letterescape}\tex{letterescape} or
    \PlaceMacro{letterbackslash}\tex{letterbackslash}.

\stopitemize

Every time \tex{hyphenatedurl} inserts a line break, it executes the
\PlaceMacro{hyphenatedurlseparator}\tex{hyphenatedurlseparator} command,
which, by default, does nothing.  But we can redefine it so that a
representative character is inserted before the line break, in a
similar way to what happens with normal words, where a hyphen is inserted
to indicate that the word continues on the next line.  For example, with:

\type{\def\hyphenatedurlseparator{\curvearrowright}}

\def\hyphenatedurlseparator{\curvearrowright}

the following particularly long web address will
have \quote{\curvearrowright} placed before the line break, as seen here:

\startnarrower\switchtobodyfont[11pt]

% Removed reference to the evil empire.
\color[blue]{\hyphenatedurl{https://wiki.contextgarden.net/Input_and_compilation/Project_and_file_management\letterhash Naming_conventions}}

\stopnarrower

In this example, to get the \quote{\letterhash} in the URL, it was
necessary to use \tex{letterhash}.

\stopdescription

\stopsubsection


% TODO: Is there a better translation of this section's title?

\startsubsection
  [
    reference=sec:createlinks,
    title=Commands that establish the link,
  ]

% TODO: "But does not write any clickable text"??  Isn't that exactly what it does?
% I rewrote this paragraph, but was I missing the point?
%
% To create links to URLs predefined with \tex{useURL}, we can use the
% command \PlaceMacro{from}\tex{from}, which is limited to establishing
% the link, but does not write any clickable text.  The default text in
% \tex{useURL} will be used as the link text.  Its syntax is:

To create links to URLs already defined with \tex{useURL}, we can use the
command \PlaceMacro{from}\tex{from}.  This inserts clickable text into the
document; the text will either be the entire URL itself, if the
two-argument form of \tex{useURL} was used, or the fourth argument, if the
four-argument form of \tex{useURL} was used.  Its syntax is:

\type{\from[Name]}

where {\tt Name} is the name previously associated with a URL using
\tex{useURL}.

To create links and associate them with a clickable text that has not been
previously defined, we have the \PlaceMacro{goto}\tex{goto} command, which
works for both internal and external links.  Its syntax is:

\type{\goto{Clickable text}[Target]}

where {\tt Clickable text} is the text to be shown in the final document
(where a mouse click will generate the jump), and {\tt Target} can be one of:

\startitemize

  \item A label from our document.  In this case \tex{goto} will generate
    the jump in a similar way as, for example, the \tex{in} or \tex{at}
    commands already examined.  But unlike those commands, no information
    associated with the label will be (automatically) displayed at the
    location where \tex{goto} is used.  For example, here is a link to
    the \quotation{Internal References} section, generated
    with \type{\goto{xyzzy}[sec:references]}: \goto{xyzzy}[sec:references].
    (The link name \quotation{xyzzy} is horribly meaningless, but by using
    that name you can be quite certain the link text was not retrieved from
    the place where the reference was originally created.)
    % TODO the verbatim \goto example about should really have
    % "\about[sec:references]" instead of "Internal References", in case
    % someone changes the name some time.  But I didn't immediately figure
    % out how to convince \type{} to do something like this:
    % \type[escape=\ABC]{this is a \ABC\curvearrowright}
    % The wiki entries for \type, \setuptyping, ..., are very unclear on
    % this point.

  \item A URL itself.  In this case it must be expressly indicated that
    it is a URL by writing the command as follows:

    \type{\goto{Clickable text}[url(URL)]}

    where URL, in turn, can be the name previously associated with a URL by
    means of \tex{useURL}, or the URL itself, in which case, when writing
    the URL, we must ensure that the \ConTeXt's reserved characters are
    written correctly in \ConTeXt.  Writing the URL according to \ConTeXt\
    rules will not affect the functionality of the link.

    For example, having already associated a URL with {\tt WhatIsCTX}, we
    can write \type{\goto{What is ConTeXt, anyway?}[url(WhatIsCTX)]}, which
    produces this link: \goto{What is ConTeXt, anyway?}[url(WhatIsCTX)].
    As a second example, the \tex{goto} command can create
    a \MyKey{mailto:} link; assuming your PDF reader supports mail links and
    your system is configured for mail links, the
    command \type{\goto{Email}[url(mailto:spam@example.com)]} will insert a
    link with text \quotation{Email} into your document, and if the link is
    clicked the PDF reader will request your mail program to start
    composing an email to the recipient {\tt spam@example.com}.

\stopitemize

\stopsubsection

\stopsection


\startsection
    [title=Creating bookmarks in the final PDF]

PDF files can have an internal bookmark list that allows the reader to see
these bookmarks in a special window of the PDF viewer program, and also to
move through the document by simply clicking on one of these bookmarks.

By default, \ConTeXt\ does not give the output PDF a list of bookmarks.
However, getting a list of bookmarks corresponding to sections in the
document is as simple as including
the \PlaceMacro{placebookmarks}\tex{placebookmarks} command, whose syntax
is:

\type{\placebookmarks[List of sections]}

where {\tt List of sections} is a comma-separated list of the section
levels that should appear in the list of contents.  For example, the command
\starttyping
\placebookmarks[chapter,section,subsection]
\stoptyping
will put a bookmark in the bookmark list for every chapter, section ans
subsection.


Keep the following observations in mind regarding this command:

\startitemize

% TODO: JD's tests indicate that it has to come after \setupinteraction (at
% least in March 2026).  Why would this be?
%%  \item According to my tests \tex{placebookmarks} does not work if it is
%%    in the preamble of the document. But, within the body of the document
%%    (between \tex{starttext} and \tex{stoptext}, or between
%%    \tex{startproduct} and \tex{stopproduct}), it doesn't matter where you
%%    place it: the bookmark list will also include the sections or
%%    subsections prior to the command. However, I believe that the most
%%    reasonable thing for a source file to be understood properly, is to
%%    place the command at the beginning.

  \item \tex{placebookmarks} must come after \tex{setupinteraction}.
  When using one or more environment files, it is probably best to
  place \tex{placebookmarks} shortly (or immediately) after
  the \tex{setupinteraction} command.

  \item Section types defined by the user (with \tex{definehead}) are not
    always located in the right place in the bookmark list.  It is
    preferable to exclude them.

  \item If the section title in any section includes an endnote or
    footnote, the text of the footnote shall be considered part of the
    bookmark.

  \item As an argument, instead of a list of sections, we can simply
    indicate the symbolic word \MyKey{all} which, as its name indicates,
    will include all the sections; however, according to my tests, this
    word, in addition to what are certainly sections, includes texts placed
    there with some non-sectioning commands, so the resulting list is
    somewhat unpredictable.

\stopitemize

Not all PDF viewer programs allow us to view bookmarks; and many that do,
do not have this feature activated by default.  Therefore, to check the
result of this function we must make sure that our PDF reader program
supports this function and has it enabled.  Acrobat, for example, does not
display bookmarks by default, although there is a button on its toolbar to
open the bookmarks display.

Finally: by default, and depending on the PDF reader, the bookmark list may
show only the top-level bookmarks; in this case, to see the subsidiary
bookmarks, a reader would have to click on some graphical element to show
bookmarks at the next level down (and so on).  If you would like more
levels to be visible by default, you can add a second argument
to \tex{placebookmarks}, indicating which levels of bookmarks will be
visible by default.  For example, the following command tells \ConTeXt\ to
include four levels of sections in the bookmark list and to make the top
three levels visible by default:

\starttyping
\placebookmarks[chapter,section,subsection,subsubsection]
               [chapter,section,subsection]
\stoptyping

\stopsection

\stopchapter

\stopcomponent

%%% Local Variables:
%%% mode: ConTeXt
%%% eval: (auto-fill-mode 1)
%%% coding: utf-8-unix
%%% TeX-master: "../introCTX_eng.tex"
%%% End:
%%% vim:set filetype=context tw=72 : %%%
